\section{Całka Perrona}
Całka Perrona zastępuje sumy całkowe porównywaniem funkcji, których pochodne Diniego leżą po obu stronach funkcji całkowanej. Daje inną charakterystykę tej samej skalarnej klasy co całka HK.
\begin{definicja}[Pochodne Diniego i funkcje porównawcze]\label{it:def:perron}
Zakładamy $a<b$; na przedziale jednopunktowym przyjmujemy całkę równą zeru.
Dla ciągłej $H:[a,b]\to\R$ definiujemy
\[
\Dn H(x)=\liminf_{\substack{y\to x\\y\ne x}}\frac{H(y)-H(x)}{y-x},\qquad
\Du H(x)=\limsup_{\substack{y\to x\\y\ne x}}\frac{H(y)-H(x)}{y-x}.
\]
Na końcach przedziału granice są jednostronne; dopuszczamy wartości nieskończone. Dla funkcji $f:[a,b]\to\R$ o skończonych wartościach \emph{funkcją górną} nazywamy ciągłą funkcję $M$ z $M(a)=0$ i $\Dn M(x)\ge f(x)$ dla wszystkich $x$; \emph{funkcją dolną} nazywamy ciągłą $m$ z $m(a)=0$ i $\Du m(x)\le f(x)$ dla wszystkich $x$. Całki górną i dolną Perrona definiujemy przez
\[
\overline{(\mathrm{Pe})\!\int_a^b}f=\inf_M M(b),\qquad
\underline{(\mathrm{Pe})\!\int_a^b}f=\sup_m m(b).
\]
Jeśli obie rodziny funkcji porównawczych są niepuste, a te kresy są równe i skończone, nazywamy $f$ całkowalną w sensie Perrona i oznaczamy tę wspólną wartość przez $(\mathrm{Pe})\int_a^b f$.
\end{definicja}
\begin{lemat}[Porządek funkcji porównawczych]\label{it:lem:perron-order}
Jeżeli $H:[a,b]\to\R$ jest ciągła i $\Dn H(x)\ge0$ dla każdego $x$, to $H$ jest niemalejąca. Dla każdej funkcji górnej $M$ i dolnej $m$ funkcja $M-m$ jest niemalejąca, a $m(b)\le M(b)$.
\end{lemat}
\begin{proof}
Dla $\eta>0$ funkcja $G(x)=H(x)+\eta x$ ma $\Dn G\ge\eta$. Gdyby $G(v)<G(u)$ przy $u<v$, jej maksimum na $[u,v]$ byłoby osiągnięte w pewnym $c<v$. W dostatecznie bliskich punktach po prawej od $c$ ilorazy różnicowe są dodatnie, co przeczy maksymalności. Tak więc $G$ jest niemalejąca; przechodząc z $\eta\downarrow0$, otrzymujemy tezę dla $H$. Z definicji ograniczeń Diniego wynika $\Dn(M-m)\ge0$, gdyż dla dostatecznie bliskich $y$ ilorazy różnicowe $M$ i $m$ są odpowiednio $\ge f(x)-\eta$ i $\le f(x)+\eta$. Zastosowanie pierwszej części i $M(a)=m(a)=0$ kończy dowód.
\end{proof}
\begin{lemat}[Mała ciągła korekta na zbiorze wyjątkowym]\label{it:lem:perron-null-correction}
Niech $F\in\AC G_*[a,b]$, $N\subset[a,b]$ będzie zbiorem
miary Lebesgue'a zero, a $f:[a,b]\to\R$ ma skończone wartości.
Dla każdego $\varepsilon>0$ istnieje ciągła niemalejąca funkcja
$H$ z $H(a)=0$, $H(b)<\varepsilon$ taka, że dla każdego $t\in N$
i wszystkich dostatecznie bliskich $y\in[a,b]$ zachodzi
\[
 |H(y)-H(t)|\ge |F(y)-F(t)|+|f(t)|\,|y-t|.
\]
\end{lemat}
\begin{proof}
Niech $[a,b]=\bigcup_{n\ge1}E_n$ będzie domkniętym pokryciem
z definicji $\AC G_*$; puste zbiory pomijamy. Najpierw dla
ustalonego $E=E_n$ konstruujemy skończoną dodatnią miarę
borelowską $\nu$ bez atomów, taką że $\nu(N\cap E)=0$ oraz
\[
 |F(y)-F(t)|\le\nu([\min(t,y),\max(t,y)])
 \qquad(t\in E,\ y\in[a,b]).
\]
Z lematu \ref{it:lem:acstar-affine-extension} wybieramy
bezwzględnie ciągłe przedłużenie $G$ funkcji $F|_E$, afiniczne
w lukach wewnętrznych i stałe w lukach brzegowych.
Miara $|G'|\lambda_1$ kontroluje moduł każdego przyrostu $G$,
na mocy twierdzenia \ref{it:thm:ac-lebesgue-primitive}.
Połóżmy $R=|F-G|$. Funkcja $R$ jest ciągła i znika na $E$.

Dla domknięcia $[l,r]$ luki wewnętrznej definiujemy na nim
\[
 L(x)=\max_{l\le z\le x}R(z),\qquad
 B(x)=\max_{x\le z\le r}R(z).
\]
Funkcje $L$ i $-B$ są niemalejące i ciągłe. Ciągłość wynika
z jednostajnej ciągłości $R$: przy przesunięciu końca
przedziału o $h$ zmiana maksimum nie przekracza $\omega_R(|h|)$.
Ponieważ $R(l)=R(r)=0$, funkcja
$K(x)=L(x)+B(l)-B(x)$ jest ciągła, niemalejąca, $K(l)=0$,
$K(r)=2\max_{[l,r]}R$ oraz
\[
 K(y)-K(l)\ge R(y),\qquad K(r)-K(y)\ge R(y).
\]
Przedłużamy $K$ przez stałe poza $[l,r]$ i oznaczamy przez
$\sigma_{[l,r]}$ jej miarę Lebesgue'a--Stieltjesa (bez atomu
w $a$). Jest bezatomowa i skupiona w otwartej luce.
W luce brzegowej z prawym końcem $r\in E$ stosujemy tylko
$K(x)=B(l)-B(x)$, gdzie $l=a$; w luce z lewym końcem $l\in E$
stosujemy tylko $K=L$, gdzie $r=b$. Zachowujemy odpowiednie
oszacowanie przyrostu od końca należącego do $E$.
Przedłużenia również są ciągłe i ich miary nie obciążają $E$.

Na każdej luce $\max R\le\osc_{[l,r]}F$: wartości $G$ leżą
między wartościami $F$ na jej końcach (w luce brzegowej
$G$ jest stała i równa wartości $F$ na końcu w $E$).
Z dowodu lematu \ref{it:lem:acstar-null-variation} suma
oscylacji na domknięciach wszystkich luk jest skończona.
Zatem miara
\[
 \nu=|G'|\lambda_1+\sum_J\sigma_J
\]
jest skończona, bezatomowa i $\nu(N\cap E)=0$; składnik
bezwzględnie ciągły nie obciąża zbiorów zerowych, a pozostałe
składniki nie obciążają $E$. Dla $t\in E$ mamy
$|F(y)-F(t)|\le |G(y)-G(t)|+R(y)$. Gdy $y\in E$,
drugi składnik znika. Gdy $y$ leży w luce, jej koniec
należący do $E$ leży pomiędzy $t$ i $y$; odpowiednia
nierówność dla $K$ kontroluje $R(y)$ przez miarę
$\sigma_J$ od tego końca do $y$. Otrzymujemy żądane
oszacowanie przez $\nu$.

Wracamy do pokrycia $E_n$ i odpowiadających mu miar $\nu_n$.
Z regularności skończonej miary borelowskiej na zwartym
przedziale, po jej uzupełnieniu, wybieramy względnie
otwarte $O_n\supset N\cap E_n$ takie, że
$\nu_n(O_n)<\varepsilon2^{-n-2}$. Nie wymaga to
borelowskości $N$: część $N\cap E_n$ leży w borelowskim
zbiorze zerowym dla $\nu_n$, otrzymanym przez przecięcie
borelowskiego zbioru $\lambda_1$-zerowego zawierającego $N$
ze zbiorem $E_n$.
Definiujemy $H_n(x)=\nu_n(O_n\cap[a,x])$.
Jest to funkcja ciągła i niemalejąca, gdyż miara jest
bezatomowa, a $H_n(a)=0$. Dla $t\in N\cap E_n$ i $y$
dostatecznie bliskiego $t$ cały odcinek między nimi leży
w $O_n$. Wtedy
$|H_n(y)-H_n(t)|\ge |F(y)-F(t)|$.

Pozostaje składnik zależny od $f$. Dla $k\ge1$ niech
$N_k=\{t\in N:k-1\le |f(t)|<k\}$. Wybieramy względnie
otwarte $U_k\supset N_k$ z
$k\lambda_1(U_k)<\varepsilon2^{-k-2}$ i kładziemy
\[
 H_0(x)=\sum_{k\ge1} k\lambda_1(U_k\cap[a,x]).
\]
Szereg jest jednostajnie zbieżny przez oszacowanie
każdego składnika; jego suma jest ciągła, niemalejąca
i zeruje się w $a$. W punkcie $t\in N_k$ odpowiedni
składnik ma lokalnie nachylenie $k$, więc
$|H_0(y)-H_0(t)|\ge |f(t)|\,|y-t|$ dla bliskich $y$.

Wreszcie $H=H_0+\sum_{n\ge1}H_n$. Także ten szereg jest
jednostajnie zbieżny, a $H(b)\le\varepsilon/2<\varepsilon$.
Wszystkie składniki są niemalejące, więc dla dowolnych
$t,y$ ich przyrosty mają ten sam znak. Wybierając $n$
z $t\in E_n$ i $k$ z $t\in N_k$, dodajemy oba lokalne
oszacowania, co dowodzi tezy.
\end{proof}

\begin{twierdzenie}[Perron i HK]\label{it:thm:perron-hk}
Funkcje całkowalne w sensie Perrona i HK tworzą tę samą klasę, a wartości całek są równe.
\end{twierdzenie}
\begin{proof}
Dowodzimy implikacji Perron $\Rightarrow$ HK. Niech $I$ będzie całką Perrona. Dla $\varepsilon>0$ wybieramy funkcje porównawcze $M,m$ z $M(b)-m(b)<\varepsilon/2$ i kładziemy $\eta=\varepsilon/(2(b-a))$. Z definicji granic Diniego dla każdego $t$ istnieje $\delta(t)>0$ taka, że dla $0<|y-t|<\delta(t)$ zachodzą nierówności
\[
\frac{M(y)-M(t)}{y-t}\ge f(t)-\eta,\qquad
\frac{m(y)-m(t)}{y-t}\le f(t)+\eta.
\]
Dla przedziału $[u,v]$ z $u\le t\le v$ zawartego w otoczeniu cechowania mnożymy pierwszą nierówność przez $v-t>0$ dla $y=v>t$, a dla $y=u<t$ przez $u-t<0$, odwracając jej znak. Dostajemy $M(v)-M(t)\ge(f(t)-\eta)(v-t)$ i $M(t)-M(u)\ge(f(t)-\eta)(t-u)$. Gdy $u=t$ lub $v=t$, odpowiednia nierówność jest równością $0=0$ i nie dzielimy przez zero. Dodanie daje $M(v)-M(u)\ge(f(t)-\eta)(v-u)$. Z drugiej nierówności analogicznie wynika $m(v)-m(u)\le(f(t)+\eta)(v-u)$. Zatem
\[
m(v)-m(u)-\eta(v-u)\le f(t)(v-u)\le M(v)-M(u)+\eta(v-u).
\]
Sumujemy po komórkach dowolnego $\delta$-drobnego podziału, uzyskując $m(b)-\eta(b-a)\le S(f,\mathcal P)\le M(b)+\eta(b-a)$. Ponieważ $m(b)\le I\le M(b)$ i $M(b)-m(b)<\varepsilon/2$, każda suma różni się od $I$ o mniej niż $\varepsilon$. Jest to całkowalność HK.

Dla implikacji HK $\Rightarrow$ Perron niech $f\in\HK[a,b]$
i $F(x)=\int_a^x f$. Z twierdzeń
\ref{it:thm:kh-primitive-derivative} i
\ref{it:thm:kh-acgstar-characterization} wiemy, że
$F\in\AC G_*[a,b]$ i $F'=f$ poza zbiorem zerowym $N$;
dołączamy do $N$ końce $a,b$. Dla $\varepsilon>0$
wybieramy korektę $H$ z lematu \ref{it:lem:perron-null-correction}
i definiujemy $M=F+H$, $m=F-H$. Obie funkcje są ciągłe
i zerują się w $a$. Poza $N$ ilorazy różnicowe $H$ są
nieujemne, dlatego $\Dn M\ge F'=f$ i $\Du m\le F'=f$.
Dla $t\in N$ oraz dostatecznie bliskiego $y\ne t$,
monotoniczność $H$ daje
\[
 \frac{H(y)-H(t)}{y-t}
 =\frac{|H(y)-H(t)|}{|y-t|}
 \ge \left|\frac{F(y)-F(t)}{y-t}\right|+|f(t)|.
\]
Zatem ilorazy różnicowe $M$ są nie mniejsze od $f(t)$,
a ilorazy $m$ nie większe od $f(t)$, także przy
zbliżaniu jednostronnym na końcach. Są to funkcje
porównawcze w sensie definicji. Ich wartości końcowe
wynoszą $F(b)\pm H(b)$. Lemat o porządku oraz
$H(b)<\varepsilon$ ściskają kresy Perrona między
$F(b)-\varepsilon$ i $F(b)+\varepsilon$. Dowolność
$\varepsilon$ daje całkowalność Perrona i wartość $F(b)$.
Dla $a=b$ obie całki przyjmujemy równe zeru.
\end{proof}
\begin{stwierdzenie}[Odtwarzanie pochodnej]
Jeżeli $F:[a,b]\to\R$ jest ciągła i różniczkowalna na $(a,b)$, a $f:[a,b]\to\R$ jest równa $F'$ we wnętrzu i ma dowolne skończone wartości na końcach, to $f$ jest całkowalna w sensie Perrona i $(\mathrm{Pe})\int_a^b f=F(b)-F(a)$.
\end{stwierdzenie}
\begin{proof}
Z twierdzenia o całkowalności pochodnych w rozdziale o całce HK wynika $f\in\HK[a,b]$ oraz $(\HK)\int_a^b f=F(b)-F(a)$, również bez istnienia skończonych pochodnych jednostronnych. Twierdzenie \ref{it:thm:perron-hk} daje tezę. Gdy pochodne jednostronne istnieją i są skończone oraz $f$ przyjmuje ich wartości na końcach, można także bezpośrednio zastosować funkcję $F-F(a)$ jako funkcję górną i dolną; lemat \ref{it:lem:perron-order} ściska obie całki do jej przyrostu.
\end{proof}
\begin{przyklad}[Funkcja charakterystyczna liczb wymiernych]
Niech $f=\ind{\Q\cap[0,1]}$. Dla $\varepsilon>0$ wybieramy otwarty zbiór $O\supset\Q\cap[0,1]$ o mierze Lebesgue'a mniejszej od $\varepsilon$ i definiujemy $M(x)=\lambda(O\cap[0,x])$, $m(x)=0$. Funkcja $M$ jest ciągła, niemalejąca i ma pochodną równą $1$ w każdym punkcie wymiernym, ponieważ punkt ten należy do wnętrza $O$. W pozostałych punktach $\Dn M\ge0=f$, a $\Du m=0\le f$ wszędzie. Rodziny funkcji porównawczych są niepuste. Z lematu \ref{it:lem:perron-order} wynika, że kres dolny wartości funkcji górnych nie jest mniejszy od kresu górnego wartości funkcji dolnych. Ponieważ $m(1)=0$ i $M(1)<\varepsilon$, oba kresy należą do $[0,\varepsilon)$. Dowolność $\varepsilon>0$ daje ich równość i skończoność, zatem całka Perrona istnieje i wynosi zero.
\end{przyklad}
\begin{uwaga}[Dlaczego warunek punktowy jest ważny]
Nie wolno bezkrytycznie zastępować warunków na pochodne funkcji porównawczych nierównościami prawie wszędzie. Funkcja Cantora i funkcja zerowa mają pochodną równą zero prawie wszędzie, ale różne przyrosty na $[0,1]$. Potrzebna jest kontrola zbiorów wyjątkowych.
\end{uwaga}

% ============================================================
% UZUPELNIENIE MATEMATYCZNE: ROZDZIAL 5
% ============================================================
\subsection*{Konstrukcja funkcji porównawczych}
\begin{lemat}[Ścisłe funkcje górne i dolne]\label{lem:perron-scisle}
Niech $f:[a,b]\to\R$ będzie funkcją Perrona-całkowalną, a $I=(\mathrm{Pe})\int_a^b f$. Dla każdego $\varepsilon>0$ istnieją funkcja górna $M$ i funkcja dolna $m$ takie, że
$0\le M(b)-I<\varepsilon$, $0\le I-m(b)<\varepsilon$ i $0\le M(x)-m(x)$ dla wszystkich $x\in[a,b]$.
\end{lemat}
\begin{proof}
Kresy definiujące całki górną i dolną Perrona są równe $I$, zatem z ich definicji można wybrać $M,m$ o wymaganych wartościach końcowych. Funkcja $M-m$ jest niemalejąca na mocy lematu o porządku z początku rozdziału i znika w $a$, więc pozostaje nieujemna na całym przedziale.
\end{proof}
\begin{twierdzenie}[Bezpośrednia konstrukcja dla funkcji z kontrolowanym zbiorem wyjątkowym]\label{tw:perron-wyjatki}
Niech $E\subset(a,b)$ będzie zwartym zbiorem miary Lebesgue'a zero, a $F\in C[a,b]$ będzie różniczkowalna w każdym punkcie $[a,b]\setminus E$ (jednostronnie na końcach). Załóżmy, że $F$ jest lipschitzowska ze stałą $L$ na pewnym otwartym otoczeniu $U$ zbioru $E$, a $f:[a,b]\to\R$ jest funkcją o skończonych wartościach, równą $F'$ poza $E$ oraz ograniczoną przez $K$ na $E$. Wtedy $f$ jest Perrona-całkowalna i
$(\mathrm{Pe})\int_a^b f=F(b)-F(a)$.
\end{twierdzenie}
\begin{proof}
Ustalmy $\varepsilon>0$ i $B>L+K+1$. Wybieramy zbiór otwarty $O$ taki, że $E\subset O$, $\overline O\subset U$ oraz $B\lambda(O)<\varepsilon$; jest to możliwe dzięki zwartości $E$ i $\lambda(E)=0$. Definiujemy $H(x)=B\lambda(O\cap[a,x])$. Funkcja $H$ jest ciągła, niemalejąca, $H(a)=0$ i $H(b)<\varepsilon$. W każdym punkcie $x\in E$ funkcja $H$ jest lokalnie afiniczna o nachyleniu $B$, a nachylenia siecznych $F$ mają moduły nie większe niż $L$. Stąd $\Dn(F+H)(x)\ge B-L>f(x)$ oraz $\Du(F-H)(x)\le L-B<f(x)$.

W punktach $x\notin E$ funkcja $F$ jest różniczkowalna; monotoniczność $H$ daje $\Dn(F+H)(x)\ge F'(x)=f(x)$ i $\Du(F-H)(x)\le F'(x)=f(x)$. Funkcje $M=F-F(a)+H$ i $m=F-F(a)-H$ są więc odpowiednio górna i dolna. Ich wartości końcowe różnią się o $2H(b)<2\varepsilon$. Z dowolności $\varepsilon$ i z definicji całki Perrona wynika teza.
\end{proof}
\begin{uwaga}[Granica argumentu]
W poprzednim dowodzie kontrola lipschitzowska na otoczeniu $E$ była potrzebna, aby wybrać \emph{jedno} skończone nachylenie $B$ dominujące wszystkie sieczne w punktach wyjątkowych. Dla dowolnej pierwotnej HK taka kontrola nie musi istnieć. Ogólna konstrukcja z lematu \ref{it:lem:perron-null-correction} zastępuje jednolite ograniczenie lipschitzowskie miarami kontrolującymi przyrosty na zbiorach pokrycia $\AC G_*$; sama tożsamość $F'=f$ prawie wszędzie nie wystarcza.
\end{uwaga}
\begin{wniosek}[Funkcje kawałkami gładkie]\label{wn:perron-kawalkami}
Jeśli $F\in C[a,b]$ i każde ograniczenie $F|_{[x_{j-1},x_j]}$ jest klasy $C^1$ względem skończonego podziału $a=x_0<\dots<x_m=b$, to dowolna funkcja $f$ równa $F'$ poza punktami podziału i mająca w tych punktach skończone wartości jest Perrona-całkowalna, a jej całka wynosi $F(b)-F(a)$.
\end{wniosek}
\begin{proof}
Niech $L=\max_{1\le j\le m}\sup_{[x_{j-1},x_j]}|F'_j|$, gdzie $F'_j$ jest pochodną odpowiedniego ograniczenia (na końcach jednostronną). Dla $u<v$ wstawiamy między nie wszystkie punkty podziału. Twierdzenie Lagrange'a na każdej części oraz nierówność trójkąta dają $|F(v)-F(u)|\le L(v-u)$, więc $F$ jest globalnie lipschitzowska. Na skończonym zbiorze $E=\{x_1,\ldots,x_{m-1}\}$ funkcja $f$ jest ograniczona. Najpierw wybieramy wartości końcowe $f(a),f(b)$ równe pochodnym jednostronnym; twierdzenie \ref{tw:perron-wyjatki} daje tezę (także przy pustym $E$). Dowolna zmiana skończonych wartości $f$ w punktach końcowych nie zmienia całki, na mocy twierdzenia \ref{it:thm:perron-hk} i lematu \ref{it:lem:kh-null-modification}.
\end{proof}
\begin{przyklad}[Wartość bezwzględna]
Dla $F(x)=|x|$ na $[-1,2]$ niech $f(x)=-1$ dla $x<0$, $f(x)=1$ dla $x>0$, a $f(0)=17$. Wówczas $f$ jest Perrona-całkowalna oraz $(\mathrm{Pe})\int_{-1}^{2}f=F(2)-F(-1)=1$. Dowód nie wymaga nadawania sztucznej pochodnej $F'(0)$.
\end{przyklad}
\subsection*{Porównanie konstrukcji Perrona i HK}
\begin{stwierdzenie}[Kontrola sum z funkcji porównawczych]\label{st:perron-sumy}
Niech $f:[a,b]\to\R$ ma skończone wartości, $M,m$ będą ciągłymi funkcjami porównawczymi z definicji Perrona, a $\eta>0$. Istnieje cechowanie $\delta$ takie, że każdy $\delta$-drobny podział znaczony $\mathcal P$ spełnia
\[
 m(b)-\eta(b-a)\le S(f,\mathcal P)\le M(b)+\eta(b-a).
\]
\end{stwierdzenie}
\begin{proof}
W punkcie $t$ wybieramy promień cechowania tak, aby ilorazy różnicowe funkcji $M$ były większe od $f(t)-\eta$, a ilorazy różnicowe funkcji $m$ mniejsze od $f(t)+\eta$, niezależnie od kierunku zbliżania się do $t$. Dla każdego $u\le t\le v$ z tej nierówności, po uwzględnieniu znaków mianowników dla $u<t$ i $v>t$, otrzymujemy
$m(v)-m(u)-\eta(v-u)\le f(t)(v-u)\le M(v)-M(u)+\eta(v-u)$. Sumowanie po komórkach kończy dowód.
\end{proof}
W ten sposób implikacja Perron $\Rightarrow$ HK jest konstruktywna: para funkcji $M,m$ bliska w punkcie $b$ daje cechowanie kontrolujące \emph{każdą} sumę całkową. Kierunek odwrotny udowodniono w twierdzeniu \ref{it:thm:perron-hk}, wykorzystując charakteryzację pierwotnych HK i małą ciągłą korektę na zbiorze wyjątkowym.

